\section*{Overview}

Greedy solutions are common in competitive programming, but the hard part is often not spotting the greedy rule. It is
proving that the local choice does not paint the solution into a corner. The exchange argument is the proof shape I use
most often for that.

\section*{When to Use It}

Reach for an exchange argument when:

\begin{itemize}
  \item the algorithm repeatedly chooses the ``best'' next item under some ordering,
  \item feasible solutions can be compared by replacing one chosen item with another,
  \item you suspect an optimal solution can be transformed to agree with the greedy choice step by step.
\end{itemize}

\section*{Core Idea}

Take an arbitrary optimal solution. If it does not use the greedy first choice, modify it so that it does, without
making the solution worse. Once the first choice is forced, recurse on the remaining subproblem.

\section*{Key Insight}

The exchange does not have to preserve the \emph{exact} same structure. It only has to preserve feasibility and total
quality. This is why greedy proofs often feel easier after the right invariant is named: earliest finish time,
smallest weight, most urgent deadline, and so on.

\section*{Worked Problem}

\paragraph{Problem.}
Given intervals \([l_i, r_i)\), choose as many pairwise non-overlapping intervals as possible.

\paragraph{Greedy rule.}
Sort by increasing right endpoint and always take the first interval that starts after the previously chosen one ends.

\paragraph{Why the rule is safe.}
Let \(g\) be the interval with the smallest ending time. Consider any optimal solution and let \(o\) be its first chosen
interval. Since \(r_g \le r_o\), replacing \(o\) with \(g\) cannot block any later interval that used to fit after
\(o\). So there exists an optimal solution that starts with \(g\).

\section*{Problem Pattern}

This proof shape also appears in:

\begin{itemize}
  \item interval scheduling,
  \item Huffman-style merging arguments,
  \item choosing edges in MST proofs,
  \item deadline scheduling and certain matroid-style problems.
\end{itemize}

\section*{Correctness Intuition}

The exchange argument is constructive: it tells you how to repair any optimal solution that disagrees with the greedy
choice. After enough repairs, you get an optimal solution that matches the greedy algorithm entirely.

\section*{Complexity Analysis}

The proof technique itself has no complexity. In the interval-scheduling example, sorting dominates and the algorithm
runs in \(O(n \log n)\).

\section*{Implementation}

The code below is the canonical interval-scheduling example. The algorithm is short, but the note matters because the
proof pattern transfers to many other greedy routines.

\section*{Common Pitfalls}

\begin{itemize}
  \item Confusing a convincing example with an actual exchange proof.
  \item Replacing one object with another without checking that all future feasibility constraints remain valid.
  \item Using an exchange argument when the problem lacks enough structure; sometimes DP is the real answer.
\end{itemize}

\section*{Variants / Extensions}

\begin{itemize}
  \item Stay-ahead proofs, where the greedy partial solution is never behind any optimal partial solution.
  \item Cut and cycle properties for graph greedy algorithms.
  \item Matroid arguments, which formalize when greedy works globally.
\end{itemize}

\section*{Practice Problems}

\begin{itemize}
  \item Select the maximum number of non-overlapping intervals.
  \item Prove the correctness of Kruskal's algorithm.
  \item Scheduling tasks with a greedy order and a swap-based proof.
\end{itemize}

\section*{References}

\begin{itemize}
  \item \href{https://usaco.guide/silver/greedy-sorting}{USACO Guide: Greedy Algorithms with Sorting}.
  \item \href{https://codeforces.com/catalog}{Codeforces Catalog}.
  \item \href{https://usaco.guide/CPH.pdf}{Competitive Programmer's Handbook}.
\end{itemize}
